% !TeX root = ../../main.tex
\section{Introduction}
\label{sec:introduction}
Let $g,n$ be nonnegative integers with $2g-2+n>0$.  Denote by
$\MM_{g,n}$ the proper Deligne--Mumford stack over $\mathbb C$ of
stable curves of genus $g$ with $n$ marked points, and write $\MM_g$
when $n=0$ and $g\geq2$.  We work in its rational Chow ring
$A^*(\MM_{g,n})_{\mathbb Q}$, where degree means complex codimension.

The Hodge bundle $\EE$ has rank $g$ and fiber $H^0(C,\omega_C)$ at
$[C;p_1,\ldots,p_n]$, where $\omega_C$ is the dualizing sheaf of the
curve $C$.  Its Chern classes are the Hodge classes
\[\lambda_j=c_j(\EE)\in A^j(\MM_{g,n})_{\mathbb Q},
\qquad 1\leq j\leq g.\]
We also set $\lambda_0=1$ and $\lambda_j=0$ for every integer $j>g$.

The stack has complex dimension $3g-3+n$.  Integration selects the
component of that codimension and takes the degree of the resulting
zero-cycle, including stabilizer weights.  Scalar
coefficients may be extended when necessary.  For a positive integer
$\ell$, we study the \emph{pure Hodge integrals}
\begin{equation}
\int_{\MM_{g,n}}\lambda_{j_1}\cdots \lambda_{j_\ell},\quad j_1, \dots, j_\ell \in \{1, \dots, g\},
\label{eq:phi-1}
\end{equation}
as $g$ tends to infinity.

The integral in \eqref{eq:phi-1} vanishes unless
\[j_1+\cdots+j_\ell=3g-3+n.\]
For integers $0\leq m<n$ with $2g-2+m>0$, let
$\pi_{n,m}:\MM_{g,n}\to\MM_{g,m}$ forget the last $n-m$ markings
and stabilize.  The Hodge bundle on the source is pulled back from the
target.  When the degree condition above holds, the corresponding
product on the target has degree $3g-3+n>3g-3+m$, so it vanishes,
as does its pullback.  Hence \eqref{eq:phi-1} vanishes for $g\geq2$,
$n\geq1$ (take $m=0$), and for $g=1$, $n\geq2$ (take $m=1$).
In genus zero there are no positive-degree Hodge classes.  The only
remaining marked case is
\[\int_{\MM_{1,1}}\lambda_1=\frac{1}{24},\]
and the large-genus problem therefore reduces to the unmarked integrals
\begin{equation}
\int_{\MM_g}\lambda_{j_1}\cdots \lambda_{j_\ell},
\quad j_1, \dots, j_\ell \in \{1, \dots, g\},
\label{eq:phi-2}
\end{equation}
where $g\ge 2$.

For $g\geq2$ and an integer $h$ with $0\leq h\leq g$, set
\[A_g=\int_{\MM_g}\lambda_1^{3g-3},\quad B_g(h)=\int_{\MM_g}\lambda_{g-h} \lambda_1^{2g-3+h},
\quad B_g=B_g(0).\]
For a homogeneous polynomial $P(\lambda)$ of Hodge classes of degree
$d$, with $0\leq d\leq3g-3$, define
\[\langle P(\lambda) \rangle_g=\frac{1}{A_g}\int_{\MM_g} P(\lambda)\lambda_1^{3g-3-d}.\]
The main results of the present paper are as follows.

\begin{thm}
\label{thm:lambda-one}
For real $x\geq0$, define
\[\Cin(x)=\int_0^x \frac{1-\cos t}{t}\dd t, \qquad
N_g=\frac{(2g-3)!(3g-3)!(4\pi^2)^{2g-3}}{\sqrt{g}(24\pi^2)^{3g-3}}.\]
As $g\to \infty$, we have
\begin{equation}
A_g=\frac{e^{\Cin(2\pi)}}{\sqrt{\pi}}N_g
\left(1+\frac{9}{8}g^{-1}+\left(\frac{257}{128}+\frac{1}{8\pi^2}\right)g^{-2}+O(g^{-3})\right).
\label{eq:111}
\end{equation}
\end{thm}

\begin{thm}
\label{thm:fixed-insertions}
Fix $\ell$ and the positive indices $j_1,\ldots,j_\ell$, and put
$P(\lambda)=\lambda_{j_1}\cdots\lambda_{j_\ell}$.  For this list of
indices, define
\[F(j)=\prod_{k=1}^\ell j_k!, \quad G(j)=\sum_{k=1}^\ell j_k(j_k-1)(j_k-2).\]
Then, as $g\to\infty$, we have
\begin{equation}
\langle P(\lambda) \rangle_g=\frac{1}{F(j)}\left(1+\frac{4\pi^2-42}{9g^2}G(j)+
\frac{4\pi^2-40}{3g^3}G(j)+O(g^{-4})\right).
\label{eq:P111}
\end{equation}
\end{thm}

\begin{thm}
\label{thm:top-hodge}
Let $g_2=2g-3$, $J_n(z)$ be the $n$-th Bessel function of the first kind, $z_{0,1}$ be the first
positive zero of $J_0$, and $\rho=6z_{0,1}J_1(z_{0,1})$.
Set
\[b=-\frac{1}{4}+\frac{1}{6z_{0,1}^2}+\frac{\pi^2}{6}J_1(z_{0,1})^2.\]
As $g\to\infty$, we have
\begin{equation}
B_g=\frac{1}{J_1(z_{0,1})}\frac{2}{(2\pi)^{2g}}\frac{(g_2!)^2}{\rho^{g_2}\sqrt{2\pi g_2}}\left(
1+\frac{b}{g_2}+O(g_2^{-2})\right).
\label{eq:g111}
\end{equation}
\end{thm}

\begin{thm}
\label{thm:fixed-defect}
With $\rho$ as above, for a fixed integer $h\geq0$, we have
\begin{equation}
\frac{B_g(h)}{B_g}\sim \frac{1}{h!}\left(\frac{24g^2}{\rho^2}\right)^h,\quad g\to\infty.
\label{eq:gh11}
\end{equation}
\end{thm}

The planned organization separates two proof routes.  After the geometric
identities in Section~2, Sections~3--5 develop uniform stable-graph and
Weil--Petersson estimates for Theorems~\ref{thm:lambda-one} and
\ref{thm:fixed-insertions}.  Section~6 treats
Theorem~\ref{thm:top-hodge} through top-Hodge intersection formulas
and Bessel analysis, independently of the Weil--Petersson asymptotics.
Sections~7--8 address Theorem~\ref{thm:fixed-defect} through linear
Hodge potentials, transport equations and singularity analysis.
